\documentclass[../../syllabus_1.tex]{subfiles}
\begin{document}

\chapter{Pourcentages et statistiques}

\section*{Matières}
\begin{enumerate}[label=\textbf{\arabic*.}]
    \item Population, variable statistique, modalités.
    \item Variable qualitative et quantitative discrète.
    \item Effectifs et fréquences.
    \item Tableau de distribution.
    \item Diagrammes statistiques.
    \item Moyenne arithmétique.
    \item Étendue.
    \item Pourcentages.
    \item Repère orthonormé.
\end{enumerate}

\section{Pourcentages}
\subsection{Découverte}

\begin{exercise}
    En 2025, \SI{45}{\%} des déchets des belges ont été recyclés ou compostés.
    \begin{enumerate}[itemsep=2cm]
        \item Sur 100 kg de déchets, combien ont été recyclés ou compostés?
        \item Sachant que les belges ont généré en moyenne 699 kg de déchets par habitant en 2025, calcule la masse des déchets recyclés ou compostés par habitant.
    \end{enumerate}
    \vspace{1.7cm}
\end{exercise}

\newpage
\begin{exercise}
    Dans un tableur, on a indiqué la répartition du traitement des déchets générés en un an par habitant en Allemagne.
    \begin{center}
        \begin{tblr}{
                hlines,
                vlines,
                colspec={lCCCCC},
                rows={.7cm, m},
                column{1}={gray!40, font=\bfseries},
                row{1}={gray!40, font=\bfseries},
                cell{1}{1}={white},
                hline{1}={1}{white},
                vline{1}={1}{white}
            }
                  & Mise en décharge & Incinérés & Recyclés & Compostés & Total \\
            En \% & 1                & 35        & 47       & 17        & 100   \\
            En kg &                  &           &          &           & 581   \\
        \end{tblr}
    \end{center}
    \begin{enumerate}[itemsep=1.2cm]
        \item Que signifie le nombre 35 ?
        \item Que signifie le nombre 581 ?
        \item Complète les cases vides du tableau.
    \end{enumerate}
\end{exercise}

\begin{exercise}
    Voici les données pour d'autres pays concernant la gestion des déchets.
    \begin{center}
        \begin{tblr}{
                hlines,
                vlines,
                colspec={lCCCCC},
                rows={.7cm, m},
                column{1}={gray!40, font=\bfseries},
                row{1}={gray!40, font=\bfseries},
                cell{1}{1}={white},
                hline{1}={1}{white},
                vline{1}={1}{white}
            }
                        & Mise en décharge en \% & Incinérés en \% & Recyclés en \% & Compostés en \% & Total en kg \\
            France      & 36                     & 32              & 18             & 14              & 543         \\
            Irlande     & 62                     & 3               & 32             & 3               & 733         \\
            Royaume-Uni & 55                     & 10              & 23             & 12              & 565         \\
        \end{tblr}
    \end{center}
    Détermine le nombre de kg non recyclés ni compostés pour chaque pays.
    \begin{enumerate}[cols=2, itemsep=2cm]
        \item France
        \item Irlande
        \item Royaume-Uni
    \end{enumerate}
\end{exercise}

\newpage
\subsection{Méthodes de calcul}
Calculer $p\,\%$ d'une quantité $N$, c'est calculer :
\begin{equation*}
    \frac{p}{100} \times N
\end{equation*}

\textbf{Exemple 1 :} Sur une tablette de chocolat noir, on peut lire « $54\,\%$ de cacao ». Quelle est la masse de cacao dans une tablette de $250$ g ?
\begin{equation*}
    \frac{54}{100} \times 250 = 135 \text{ g de cacao}
\end{equation*}

\textbf{Exemple 2 :} Un ordinateur reconditionné coûte $280$ € hors TVA. Quel sera son prix TVA ($21\,\%$) comprise ?
\begin{equation*}
    280 + \frac{21}{100} \times 280 = 280 + 58{,}8 = 338{,}8 \text{ €}
\end{equation*}

\textbf{Exemple 3 :} En faisant ses courses pendant les soldes, Sonia achète un blouson qui est affiché à 40~€. Il est indiqué d'une pastille orange, ce qui signifie \og -30~\% \fg. Quel est le prix que Sonia va payer à la caisse?
\vspace{3cm}

\begin{exercise}
    Calcule.
    \begin{enumerate}[cols=3, itemsep=2cm]
        \item 36~\% de 25 km
        \item 78~\% de 12 L
        \item 25~\% d'une heure
        \item 95~\% de 750 g
        \item 25~\% de 12~€
        \item 2~\% de \num{1500}~€
        \item 50~\% de 438
        \item 20~\% de 45 L
        \item 5~\% de 48 km
    \end{enumerate}
    \vspace{2cm}
\end{exercise}

\begin{exercise}
    Dans un collège de 575 élèves, 28~\% des collégiens sont en 6\ieme. Calcule le nombre d'élèves que cela représente.
    \vspace{3cm}
\end{exercise}

\begin{exercise}
    Une citerne ayant une capacité de \num{8500} L est remplie d'eau à 60~\%.
    \begin{enumerate}[itemsep=2cm]
        \item Quelle quantité d'eau, en litres, cette citerne contient-elle?
        \item Quelle quantité d'eau, en litres, cette citerne peut-elle encore recevoir?
    \end{enumerate}
    \vspace{2cm}
\end{exercise}

\begin{exercise}
    Construis ci-dessous un rectangle de 6 cm de longueur et de 5 cm de largeur. Hachure ensuite 40~\% de la surface de ton rectangle.
    \vspace{5cm}
\end{exercise}

\begin{exercise}
    Au cours du dernier trimestre, une usine a créé \num{15200} réfrigérateurs. Le service après-vente a noté des dysfonctionnements sur 608 d'entre eux. \\
    Détermine le pourcentage d'appareils défectueux.
    \vspace{5cm}
\end{exercise}

\begin{exercise}
    Joanna veut acheter une nouvelle voiture qui coûte \num{12000}~€. Puisqu'elle paye directement, le vendeur veut bien lui accorder une réduction de 12~\% sur le prix total. Quel sera le prix à payer pour la voiture~? \\ Laisse ta démarche visible.
    \vspace{4cm}
\end{exercise}

\begin{exercise}
    Voici quatre situations de réductions ou d'augmentation de prix. Calcule le prix final en une seule étape.
    \begin{enumerate}[cols=2, itemsep=3cm]
        \item +20~\% sur \num{5500}~€
        \item -15~\% sur \num{2000}~€
        \item Un intérêt de 3~\% sur un placement de \num{1000}~€
        \item Une remise de 5~\% sur un achat de 25~€
    \end{enumerate}
    \vspace{4cm}
\end{exercise}

\begin{exercise}
    Charlotte travaille dans un magasin où elle reçoit un salaire mensuel de \num{1400}~€. Comme elle travaille bien, son patron décide de l'augmenter, et Charlotte reçoit 6~\% de plus pour son salaire mensuel. Que gagne Charlotte chaque mois~? \\ Laisse ta démarche visible.
\end{exercise}
\vfil

\newpage
\subsection{Comparer des offres}

Pour comparer deux situations, il est souvent utile de tout ramener à un pourcentage ou à une même unité de comparaison.

\begin{exercise}
    Dans l'école Prévert, il y a $637$ élèves, dont $217$ qui font du sport dans une équipe de l'école. Dans l'école Rimbaud, il y a $480$ élèves, dont $157$ qui sont dans des associations sportives.

    Quelle est l'école la plus sportive ? Justifie en t'aidant des pourcentages.
    \vspace{3cm}
\end{exercise}

\begin{exercise}
    Cyril voit deux offres pour la toute nouvelle console.
    \begin{itemize}[label=\faGamepad]
        \item Chez \og MoinCher \fg, la console coûte 350€ mais on peut recevoir 10~\% de réduction si on commande par internet.
        \item Chez \og Bapri \fg, la console coûte 400~€, mais si la commande est faite avant ce soir, le client reçoit 20~\% de réduction.
    \end{itemize}
    Dans quel magasin Cyril doit-il aller pour payer sa console au moindre prix? Justifie.
    \vspace{3cm}
\end{exercise}

\begin{exercise}
    Deux magasins de sport proposent la même paire de baskets.
    \begin{itemize}
        \item Chez « SportPlus », la paire coûte $85$ € et bénéficie d'une réduction de $15\,\%$.
        \item Chez « RunShop », la paire coûte $90$ € et bénéficie d'une réduction de $20\,\%$.
    \end{itemize}
    Dans quel magasin le client paiera-t-il le moins cher ? Justifie ta réponse en laissant apparaitre tous tes calculs.
    \vspace{3cm}
\end{exercise}

\newpage
\subsection{Pourcentages successifs}
Lorsqu'on applique deux pourcentages l'un après l'autre, on ne peut pas simplement les additionner ! Chaque pourcentage se calcule sur le montant obtenu après la première opération.

\begin{exercise}
    Yousra a placé $3\,000$ € en banque. Après $1$ an, elle reçoit un intérêt de $7\,\%$. Ensuite, les impôts arrivent et elle doit payer une taxe de $2\,\%$ sur le nouveau montant qu'elle a en banque.

    De quel montant Yousra dispose-t-elle à la fin de ces deux opérations ? Laisse ta démarche visible.
    \vspace{3cm}
\end{exercise}

\begin{exercise}
    Un magasin propose d'abord une réduction de $10\,\%$ sur un article de $80$ €, puis, la semaine suivante, une réduction supplémentaire de $20\,\%$ sur le nouveau prix.
    \begin{enumerate}[itemsep=1cm]
        \item Quel est le prix après la première réduction ?
        \item Quel est le prix final après la seconde réduction ?
        \item Ce prix final correspond-il à une réduction totale de $30\,\%$ sur le prix de départ ? Justifie.
    \end{enumerate}
    \vspace{1cm}
\end{exercise}

\begin{exercise}
    Le prix d'un vélo est de $250$ €. Le magasin applique d'abord une augmentation de $8\,\%$ liée à une pénurie de matériaux, puis, un mois plus tard, une réduction de $15\,\%$ lors d'une offre spéciale.
    \begin{enumerate}[itemsep=1cm]
        \item Quel est le prix du vélo après l'augmentation de $8\,\%$ ?
        \item Quel est le prix du vélo après la réduction de $15\,\%$ appliquée sur ce nouveau prix ?
        \item Ce prix final correspond-il à une simple réduction de $7\,\%$ ($15\,\% - 8\,\%$) sur le prix de départ ? Justifie ta réponse par le calcul.
    \end{enumerate}
\end{exercise}

\newpage
\section{Repérage d'un point dans un plan}
Pour repérer un point sur un plan, nous avons besoin de plusieurs informations. Explique oralement, dans chacun des exemples suivants, comment repérer la croix.
\begin{multicols}{2}
    Exemple 1\\\includegraphics[width=\linewidth]{chapitres/traitement_de_donnes_et_pourcentages/figures/exemple1.png}

    Exemple 2\\\includegraphics[width=\linewidth]{chapitres/traitement_de_donnes_et_pourcentages/figures/exemple2.png}

    Exemple 3\\\includegraphics[width=\linewidth]{chapitres/traitement_de_donnes_et_pourcentages/figures/exemple5.png}

    \columnbreak
    Exemple 4\\\includegraphics[rotate=-90,width=\linewidth]{chapitres/traitement_de_donnes_et_pourcentages/figures/exemple3.png}
\end{multicols}

\newpage
En mathématiques, pour repérer un point sur un plan, nous utilisons majoritairement un repère dit \textbf{orthonormé} (appartenant à la famille des repères cartésiens). Un repère orthonormé possède deux caractéristiques majeures: ses axes sont perpendiculaires et ils ont les mêmes normes. L'intersection des deux axes se nomme \textbf{l'origine du repère}.

Pour placer ou repérer un point dans un repère orthonormé \emph{(exemple 3)}, nous avons besoin des \textbf{coordonnées} de ce point. \\

Ces coordonnées sont données dans un ordre précis~: tout d'abord en donnant \textbf{l'abscisse} du point ($x$), qui est son repérage sur l'axe horizontal (axe des abscisses), et en donnant ensuite \textbf{l'ordonnée} du point ($y$), qui est son repérage sur l'axe vertical (axe des ordonnées).

\tcbsidebyside[sidebyside adapt=right,
    blanker,sidebyside gap=5mm
]{%
    Les coordonnées se notent toujours sous la forme $(x;y)$.
    \bigbreak
    Par exemple, le point $A$ a comme abscisse 4, et comme ordonnée 2. On notera alors $A(4;2)$.
    \bigbreak
    Note les coordonnées des points $B$, $C$ et $D$.
}{%
    \includegraphics{chapitres/traitement_de_donnes_et_pourcentages/figures/axes.tikz}
}

\begin{exercise}
    Note correctement les coordonnées des points suivants en dessous du repère.
    \begin{center}
        \includegraphics{chapitres/traitement_de_donnes_et_pourcentages/figures/ex_reperage1.tikz}
    \end{center}
\end{exercise}

\newpage
\begin{exercise}
    Place dans le repère suivant les points dont voici les coordonnées.

    $D\left(-2;3\right)$
    \hfill
    $A\left(3;0\right)$
    \hfill
    $N\left(2;-2\right)$
    \hfill
    $I\left(1;1\right)$
    \hfill
    $E\left(4;3\right)$
    \hfill
    $L\left(-3;-2\right)$

    \begin{center}
        \includegraphics[width=\linewidth]{chapitres/traitement_de_donnes_et_pourcentages/figures/ex_reperage2.tikz}
    \end{center}
\end{exercise}

\begin{exercise}
    Donne les coordonnées des points qui satisfont les conditions ci-dessous.
    \begin{enumerate}[itemsep=\fill]
        \item Mon abscisse vaut 4 et mon ordonnée vaut la moitié de mon abscisse.
        \item Mon abscisse est un multiple de 3. Mon ordonnée est un multiple de 4. La somme de mes coordonnées vaut 10.
        \item Mon ordonnée vaut 3 et mon abscisse vaut 6.
        \item Mon ordonnée vaut 6 et mon abscisse vaut le tiers de mon ordonnée.
        \item Mon abscisse vaut 3 et mon ordonnée vaut son carré.
    \end{enumerate}
\end{exercise}

\newpage
\begin{exercise}
    Dans le quadrillage suivant, construis un repère cartésien. Tu utiliseras comme échelle \emph{1 carré = 1 unité}. Ton axe des abscisses sera gradué de $-4$ à $4$ et ton axe des ordonnées de $-2$ à $5$.\\
    Place ensuite les points suivants.

    $I\left(-1;-1\right)$
    \hfill
    $L\left(-3;1\right)$
    \hfill
    $S\left(2;0\right)$
    \hfill
    $A\left(3;4\right)$
    \hfill
    $E\left(-4;4\right)$

    \begin{center}
        
\includegraphics[]{chapitres/traitement_de_donnes_et_pourcentages/figures/quadrillage.tikz}
    \end{center}
\end{exercise}

\newpage

\section{Lecture de graphiques, analyse de données}
Dans ce chapitre, nous rencontrerons 3 types de graphique qui permettent de visualiser des données.
\begin{enumerate}
    \item Le graphique évolutif. Ce type de graphique s'utilise pour montrer l'évolution d'une grandeur par rapport à une autre.
          \begin{center}
              \includegraphics{tikz/diagramme_evolutif_theorie.tikz}
          \end{center}
    \item Le diagramme en bâton.Ce diagramme s'utilise pour comparer des données.
          \begin{center}
              \includegraphics{tikz/diagramme_baton_theorie.tikz}
          \end{center}
    \item Le diagramme circulaire. Ce diagramme s'utilise pour représenter des données dont on connaît la répartition par rapport à un tout.
          \begin{center}
              \includegraphics{tikz/diagramme_circulaire_theorie.tikz}
          \end{center}
\end{enumerate}

\newpage

\begin{exercise}
    Un concessionnaire automobile a vendu ce mois-ci 85 véhicules de tous types. En voici un descriptif partiel.
    \begin{center}
        \begin{tblr}{
                colspec={CCCCC},
                hlines,
                vlines,
                row{1}={gray!40, font = \bfseries},
                column{1}={gray!40,font = \bfseries},
                rows={1cm, m}}
            Vendeurs & Citadines & Sportives & Routières & Totaux \\
            Paul     & 3         & 5         &           & 17     \\
            Denis    & 4         &           & 6         & 15     \\
            Henri    & 3         &           & 8         &        \\
            Steeve   &           & 4         &           & 18     \\
            Eliess   & 5         &           & 2         & 16     \\
            Totaux   &           & 31        & 30        & 85     \\
        \end{tblr}
    \end{center}
    \begin{enumerate}[itemsep=1em]
        \item Complète le tableau.
        \item Combien de voitures Henri a-t-il vendu?
        \item Combien de citadines ont été vendues dans cette concession?
        \item Quel est le vendeur qui a vendu le plus de sportives?
        \item Denis est persuadé d'avoir vendu autant de sportives que de routières. A-t-il raison?
        \item Qui est le meilleur vendeur?
    \end{enumerate}
\end{exercise}

\begin{exercise}
    Un agriculteur a réalisé le diagramme circulaire suivant illustrant l'utilisation des terres de son exploitation.

    \tcbsidebyside[sidebyside adapt=right,
        blanker,sidebyside gap=1cm
    ]{%
        \begin{enumerate}[itemsep=1cm]
            \item Quel type de culture occupe la moitié de ses terres?
            \item Quel type de culture est la moins répandue sur ses terres?
            \item Quelles sont les cultures qui occupent la même surface?
        \end{enumerate}
    }{%
        \includegraphics[width=.5\linewidth]{chapitres/traitement_de_donnes_et_pourcentages/figures/lecture_de_donnees2.tikz}
    }
\end{exercise}

\newpage
\begin{exercise}
    Observe le graphique évolutif ci-dessous illustrant le prix de location d'une planche à voile, et complète le tableau en commençant par donner un nom aux têtes de colonnes.

    \tcbsidebyside[
        sidebyside adapt=left,
        blanker,
        sidebyside gap=1cm
    ]{%
        \includegraphics[width=10cm]{chapitres/traitement_de_donnes_et_pourcentages/figures/lecture_de_donnees3.tikz}
    }{%
        \begin{tblr}{hlines, vlines, rows={1cm, m}, columns={2cm, c}, row{1}={2cm, m}}
              & \\
            1 & \\
            2 & \\
            3 & \\
            4 & \\
            5 & \\
            6 & \\
        \end{tblr}
    }
\end{exercise}

\begin{exercise}
    Le diagramme circulaire suivant montre la répartition des élèves de la même classe en fonction du nombre de sports pratiqués régulièrement.

    \includegraphics[width=\linewidth]{chapitres/traitement_de_donnes_et_pourcentages/figures/lecture_de_donnees4.tikz}

    \begin{enumerate}[cols=2]
        \item Si tu sais que 3 élèves ne pratiquent aucun sport régulièrement, détermine le nombre d'élèves de la classe.\columnbreak
        \item Combien de personnes pratiquent plus de un sport régulièrement~?
    \end{enumerate}


\end{exercise}

\newpage
\section{Construction de graphiques}
Dans une école de Wallonie, en 1\iere{} année, il y a 120 élèves. Ils ont eu le choix entre plusieurs options : il en a 30 qui sont en informatique, 25 qui suivent le latin, 40 qui font un renforcement math et le reste qui sont en activité complémentaire d'expression corporelle. Représente la répartition des élèves sur un diagramme circulaire. \par \par
\emph{Pour réaliser ce type de diagramme, nous allons devoir déterminer dans un premier temps le pourcentage que chaque option représente par rapport à l'ensemble des élèves de 1\iere{} années, puis transformer cela en amplitudes pour pouvoir le placer dans le diagramme.}

\begin{center}
    \begin{tblr}{
            hlines,
            vlines,
            colspec={CCCC},
            rows={1cm, m},
            column{1}={gray!40, font=\bfseries},
            row{1}={2cm, gray!40, font=\bfseries},
        }
        Option                & Nombre d'élèves & Pourcentage & Amplitude \\
        Informatique          &                 &             &           \\
        Latin                 &                 &             &           \\
        Renforcement math     &                 &             &           \\
        Expression corporelle &                 &             &           \\
        Total                 &                 &             &           \\
    \end{tblr}
\end{center}

\emph{Représente maintenant les répartitions des élèves dans un diagramme circulaire. Fais une légende de ton diagramme.} \\


\includegraphics[width=6cm]{chapitres/traitement_de_donnes_et_pourcentages/figures/cercle.tikz}

\newpage
\begin{exercise}
    En Belgique, sur 12 millions d'habitants, 6,6 millions parlent néerlandais comme première langue, 0,5 millions parlent allemand et 3,7 millions ont le français comme langue maternelle. Quel est pourcentage de la population qui ne parle aucune de ces 3 langues ? Représente la situation dans un diagramme circulaire. \\
    \emph{Note : retranscris d'abord tes données sous la forme d'un tableau.}

    
\includegraphics{chapitres/traitement_de_donnes_et_pourcentages/figures/cercle.tikz}
\end{exercise}
\vfill
\begin{exercise}
    Éléanore a relevé les températures dehors au cours de la journée. Elle les écrit comme suit.
    \tcbsidebyside[
        sidebyside adapt=left,
        blanker,
        sidebyside gap=1cm
    ]{%
        \begin{tblr}{
                hlines,
                vlines,
                rowsep={1em}
            }
            \SI{8}{\hour}  & \SI{10}{\degreeCelsius} \\
            \SI{10}{\hour} & \SI{11}{\degreeCelsius} \\
            \SI{12}{\hour} & \SI{12}{\degreeCelsius} \\
            \SI{14}{\hour} & \SI{17}{\degreeCelsius} \\
            \SI{16}{\hour} & \SI{15}{\degreeCelsius} \\
            \SI{18}{\hour} & \SI{11}{\degreeCelsius} \\
        \end{tblr}
    }{%
        
\includegraphics[width=\linewidth]{chapitres/traitement_de_donnes_et_pourcentages/figures/quadri2.tikz}
    }

    Trace le graphique évolutif qui correspond à la situation.
\end{exercise}

\newpage
\begin{exercise}
    Lorsque le robinet de ma terrasse reste ouvert au maximum, il remplit trois seaux de 10 litres en une minute. \\ Complète le tableau et construis le graphique.
    \begin{center}
        \begin{tblr}{
                hlines,
                vlines,
                columns={2cm, c},
                rows={1cm, m},
                column{1}={3cm, gray!40, font=\bfseries}
            }
            Temps en minutes         & 1 & 2 & 3 & 7 & 15 \\
            Quantité d'eau en litres &   &   &   &   &    \\
        \end{tblr}
    \end{center}

    Titre du graphique : \dotfill

    \begin{center}
        \includegraphics[height=15cm]{chapitres/traitement_de_donnes_et_pourcentages/figures/robinet.tikz}
    \end{center}
\end{exercise}

\newpage
\begin{exercise}
    Lors d'une enquête menée par Anais parmi les élèves de 3\ieme{} maternelle, elle a noté les couleurs préférées des enfants.
    \begin{center}
        \begin{tblr}{
                hlines,
                vlines,
                columns={1.5cm, c},
                rows={.7cm, m},
                row{1}={gray!40, font=\bfseries}
            }
            Jaune & Rouge & Rose & Vert & Bleu & Orange & Autre & Total \\
            3     & 2     & 5    & 2    & 7    & 5      & 3     &       \\
        \end{tblr}
    \end{center}
    \vfill

    \begin{enumerate}
        \item Combien d'enfants se trouvent dans la classe?
        \item Détermine le pourcentage d'enfants qui préfèrent le bleu.
        \item Construis un diagramme circulaire qui représente la situation.
              \begin{center}
                  
\includegraphics{chapitres/traitement_de_donnes_et_pourcentages/figures/cercle.tikz}
              \end{center}
        \item Construis un diagramme en bâtonnets qui représente la situation.\bigbreak
              \begin{center}
                  
\includegraphics[width=.5\linewidth]{chapitres/traitement_de_donnes_et_pourcentages/figures/quadri2.tikz}
              \end{center}
    \end{enumerate}
\end{exercise}

\newpage
\newpage
\section*{Mots-clès}
\begin{enumerate}[cols=2, label=\textbf{\arabic*.}]
    \item Population
    \item Modalité
    \item Variable qualitative
    \item Variable quantitative
    \item Effectif, effectif total
    \item Fréquence
    \item Tableau de distribution
    \item Diagramme en bâtonnets
    \item Diagramme à bandes
    \item Diagramme circulaire
    \item Moyenne arithmétique
    \item Étendue
    \item Pourcentage
    \item Repère
    \item Abscisse
    \item Ordonnée
\end{enumerate}

\section*{Attendus}
\begin{enumerate}[label=\textbf{\arabic*.}]
    \item Identifier la population, la variable, les modalités, les effectifs, les fréquences et l'étendue.
    \item Caractériser la variable étudiée (qualitative ou quantitative discrète).
    \item Identifier le diagramme donné : diagramme en bâtonnets (variable quantitative discrète), diagramme à bandes (variable qualitative) ou diagramme circulaire.
    \item Décrire le concept de moyenne arithmétique (variable discrète).
    \item À partir d'une situation en français, d'un tableau de distribution ou d'un diagramme statistique : décrire la population et la variable étudiées, les caractériser, et déterminer l'effectif total, l'effectif d'une modalité ou sa fréquence.
    \item Calculer la moyenne arithmétique d'une variable quantitative discrète.
    \item Relier entre elles différentes présentations d'une même situation (liste de données, tableau de distribution, diagrammes).
    \item Présenter une liste de données à l'aide d'un diagramme en bâtonnets ou à bandes.
    \item Interpréter une valeur obtenue en lien avec le caractère étudié et le contexte.
    \item Lire des informations présentées à partir de supports différents pour répondre à des questions.
    \item Choisir, parmi un ensemble d'informations (moyenne, effectifs, fréquence, diagrammes, tableau de distribution), l'élément pertinent pour répondre à une question posée.
    \item Calculer un pourcentage en situation contextualisée.
    \item Comparer des offres faisant intervenir des pourcentages afin d'en déterminer la plus avantageuse.
    \item Exploiter une situation faisant intervenir des pourcentages successifs.
    \item Énoncer les éléments d’un repère orthonormé en utilisant le vocabulaire adéquat (axes perpendiculaires, unités, origine).
    \item Identifier l’abscisse, l’ordonnée, les coordonnées à partir d’un point placé dans un repère ou d’un couple de nombres.
    \item Écrire, avec le symbolisme adéquat, les coordonnées d’un point.
    \item Déterminer les coordonnées d’un point placé dans un repère orthonormé.
    \item Placer, dans un repère orthonormé, un point dont les coordonnées sont données.
    \item Placer un ensemble de points qui vérifient une condition exprimant, en langage courant, une relation entre abscisse et ordonnée. Ex. : l’ordonnée est le double de l’abscisse.
\end{enumerate}

\end{document}